\documentclass[11pt,reqno]{amsart}
\usepackage[T1]{fontenc}
\usepackage{lmodern}
\usepackage{microtype}
\usepackage{amsmath,amssymb,mathtools}
\usepackage{booktabs}
\usepackage[colorlinks=true,linkcolor=blue,citecolor=blue,urlcolor=blue]{hyperref}
\newtheorem{theorem}{Theorem}[section]
\newtheorem{lemma}[theorem]{Lemma}
\newtheorem{proposition}[theorem]{Proposition}
\newtheorem{corollary}[theorem]{Corollary}
\newtheorem{conjecture}[theorem]{Conjecture}
\newtheorem*{thmQ}{Theorem Q}
\newtheorem*{thmE}{Theorem E}
\newtheorem*{thmU}{Theorem U}
\newtheorem*{thmW}{Theorem W}
\newtheorem*{thmN}{Theorem N}
\theoremstyle{definition}
\newtheorem{definition}[theorem]{Definition}
\newtheorem{remark}[theorem]{Remark}
\newtheorem{computation}[theorem]{Computation}
\newtheorem{citedthm}[theorem]{Cited theorem}
\newtheorem{hypothesis}[theorem]{Hypothesis}
\newtheorem{problem}{Problem}
\newcommand{\Z}{\mathbb Z}
\newcommand{\Q}{\mathbb Q}
\newcommand{\N}{\mathbb N}
\newcommand{\cC}{\mathcal C}
\newcommand{\cS}{\mathcal S}
\newcommand{\cQ}{\mathcal Q}
\newcommand{\cA}{\mathcal A}
\newcommand{\cK}{\mathcal K}
\newcommand{\cT}{\mathcal T}
\newcommand{\HC}{\operatorname{HC}}
\newcommand{\li}{\operatorname{li}}
\newcommand{\leg}[2]{\left(\frac{#1}{#2}\right)}
\newcommand{\EH}{\mathrm{EH}}
\title[The least certificate coordinate for Erd\H{o}s--Straus]{The least
certificate coordinate for the Erd\H{o}s--Straus equation at a prime:
a two-shift reduction, partial sieve results, and covering obstructions}
\author{Anonymous}
\date{Draft, November 2026; internal checks only, not externally refereed}
\emergencystretch=3em
\begin{document}
\begin{abstract}
For an integer $n\equiv1\pmod 4$ call a triple $(t,u,v)$ of positive
integers with $tuv-u-v=n$ a \emph{word} of $n$, and let $\cC_n$ be the set
of positive integers that are $\equiv0\pmod4$ or $\equiv-n\pmod8$.  A word
with a coordinate in $\cC_n$ yields an explicit solution of
$4/n=1/x+1/y+1/z$, and $c\in\cC_n$ is a coordinate of a word if and only if
$n+c$ or $cn+1$ has a divisor $\equiv-1\pmod c$.  We study the least such
coordinate $\HC(p)$ at primes $p$.  It equals $3$ for $p\equiv5\pmod8$
and is at most $8$ unless $p$ lies in one of Mordell's six classes modulo
$840$.  For $p\equiv1\pmod8$ we show: $\HC(p)>4$ holds exactly when $p+4$
and $4p+1$ have only prime factors $\equiv1\pmod4$; every prime
$p=5m^2+1$ with $4\mid m$ has $\HC(p)\ge7$; under the Elliott--Halberstam
conjecture at some level $1-\eta_0<1$ there are $\gg x/(\log x)^2$ primes
$p\le x$ with $\HC(p)>4$, which is the true order; and unconditionally
$\gg x/(\log x)^2$ primes $p\le x$, $p\equiv1\pmod 8$, have $p+4$ composed
of primes $\equiv1\pmod4$ and $4p+1$ free of primes $\equiv3\pmod4$ below
$x^{\delta}$ for an absolute $\delta>0$.  We show that no finite family of
fixed certificates covers the primes in a class
$p\equiv1\pmod{840}$, $p\equiv R\pmod\ell$ with $R$ a small quadratic
non-residue modulo a prime $\ell\ge11$; such classes contain no squares,
so this obstruction is not the classical one at squares.  We also bound
the reach of Dickson-type constructions of primes with large $\HC$.  A
computation gives $\HC(p)\le64$ for every prime $p\equiv1\pmod4$ below
$10^{12}$.  The Erd\H{o}s--Straus conjecture remains open; nothing here
proves that $\HC(p)$ is finite for all large $p$, nor that it is unbounded.
\end{abstract}
\maketitle

\section{Introduction}\label{sec:intro}

The Erd\H{o}s--Straus conjecture asserts that for every integer $n\ge2$
\begin{equation}\label{eq:ES}
\frac4n=\frac1x+\frac1y+\frac1z
\end{equation}
has a solution in positive integers.  It is open.  It suffices to treat
primes, and classical polynomial identities settle every prime outside
the residue classes $1,121,169,289,361,529\pmod{840}$ (Mordell
\cite{Mordell}; see \cite[p.~8]{ET}).  Finite families of such identities
cannot settle any class consisting of squares
(Yamamoto~\cite{Yamamoto}, Schinzel~\cite{Schinzel}; see
\cite[Prop.~1.6 and the discussion before Prop.~1.9]{ET}).  The equation
has been verified numerically for $n\le10^{17}$ \cite{Salez}, and a recent
preprint reports $n\le10^{18}$ \cite{MD}.  On average there are many
solutions: Elsholtz and Tao \cite[Thm.~1.1]{ET} show
$N(\log N)^2\ll\sum_{p\le N}f(p)\ll N(\log N)^2\log\log N$ for the number
$f(p)$ of solutions.

This note isolates one family of explicit solutions and studies the size
of its smallest parameter.  For $n\equiv1\pmod4$ a \emph{word} of $n$ is a
triple $(t,u,v)\in\N^3$ with
\[
tuv-u-v=n,
\]
and we put
\[
\cC_n:=\{x\in\N:\ x\equiv0\ (\mathrm{mod}\ 4)\ \text{or}\ x\equiv-n\ (\mathrm{mod}\ 8)\}.
\]
If some coordinate of a word lies in $\cC_n$, then one of five explicit
identities solves \eqref{eq:ES} (Proposition~\ref{prop:slots}).  The
coordinate $c\in\cC_n$ occurs in a word exactly when $n+c$ or $cn+1$ has
a positive divisor $d\equiv-1\pmod c$ (Lemma~\ref{lem:lines}); we call
such a pair $(c,d)$ a \emph{certificate} and define
\[
\HC(n):=\min\{c:\ (c,d)\ \text{is a certificate for}\ n\}\in\N\cup\{\infty\}.
\]
So $\HC(p)<\infty$ implies that \eqref{eq:ES} is solvable for $n=p$.  We
do not know whether the converse holds.  The coordinates divisible by
$4$ are the ``Type~A'' and ``Type~B'' solutions of L\'opez \cite{Lopez},
who conjectures that every prime has one (Remark~\ref{rem:lopez}).

$\HC(p)=3$ for $p\equiv5\pmod 8$, so the interesting primes are
$p\equiv1\pmod 8$, where $\cC_p=\{4,7,8,12,15,16,20,23,\dots\}$.  A
uniform bound $\HC(p)\le F(p)$ for all large $p$, for any function $F$,
would prove the Erd\H{o}s--Straus conjecture for all large primes.  We
prove no such bound.  Our results concern the following questions: how
small is $\HC$ on explicit classes, how often and by which mechanism can
it be large, and what can a covering argument achieve.

\subsection*{Results}
Throughout, $\cQ$ is the set of primes $\equiv3\pmod4$ and $\cS$ is the
set of positive integers all of whose prime factors are $\equiv1\pmod4$.

\begin{enumerate}
\item\emph{Basic structure (\S\ref{sec:words}).}  Every certificate divisor
$d$ for $n$ has Jacobi symbol $(d/n)=-1$ (Lemma~\ref{lem:recip}).  Hence
$\HC(m^2)=\infty$ for every odd $m$, and certificates supported on primes
dividing $2m$ never occur when $n\equiv1\pmod{8m}$.
\item\emph{Explicit classes (\S\ref{sec:classes}).}  $\HC(p)\le8$ for every
prime $p>7$, $p\equiv1\pmod 4$, outside Mordell's six classes modulo
$840$ (Proposition~\ref{prop:mordell}), and
$\#\{p\le x:\ p\equiv1\ (8),\ \HC(p)>4\}\ll x/(\log x)^2$
(Proposition~\ref{prop:upper4}).
\item\emph{The first threshold (\S\ref{sec:kill}).}  For $p\equiv1\pmod8$,
$\HC(p)>4$ if and only if $p+4\in\cS$ and $4p+1\in\cS$
(Lemma~\ref{lem:K4}).  So ``$\HC(p)>4$ for infinitely many $p$'' is the
two-shift prime problem
\begin{equation}\label{eq:P415}
\text{infinitely many primes } p\equiv1\ (\mathrm{mod}\ 8)\ \text{with}\ p+4\in\cS\ \text{and}\ 4p+1\in\cS,
\end{equation}
an analogue of the hyperbolic prime number theorem of Friedlander and
Iwaniec \cite{FI09} (Remark~\ref{rem:shell}).  We prove:
\begin{itemize}
\item \emph{Theorem Q.} Every prime $p=5m^2+1$ with $4\mid m$ has
$\HC(p)\ge7$.  Hence \eqref{eq:P415} follows from Bunyakovsky's conjecture
for the single polynomial $720k^2+1$.
\item \emph{Theorem E.} There is an absolute $\eta_0>0$ such that the
Elliott--Halberstam hypothesis $\EH(1-\eta_0)$ implies
$\#\{p\le x:\ p\equiv1\ (8),\ p+4\in\cS,\ 4p+1\in\cS\}\gg x/(\log x)^2$.
With Proposition~\ref{prop:upper4} this is the true order.
\item \emph{Theorem U} (unconditional).  There are absolute $\delta>0$ and
$c>0$ with
\[
\#\{p\le x:\ p\equiv1\ (8),\ p+4\in\cS,\ 4p+1\ \text{has no prime factor in}\ \cQ\ \text{below}\ x^{\delta}\}\ge c\,\frac{x}{(\log x)^2}
\]
for large $x$.  For these $p$, $4p+1$ has a bounded number of prime
factors in $\cQ$; \eqref{eq:P415} needs that number to be $0$.
\end{itemize}
\item\emph{Dickson-type constructions (\S\ref{sec:dickson}).}  Under
Dickson's conjecture $\HC$ is unbounded (Proposition~\ref{prop:dickson}).
Theorem~W shows that a fixed polynomial family of this kind can certify
$\HC(p)>X$ only for $X<(4+o(1))\log p$, except possibly at its least
member.
\item\emph{Coverings (\S\ref{sec:cover}).}  \emph{Theorem N}: let
$\ell\ge11$ be prime and $R$ a quadratic non-residue modulo $\ell$ with
$1\le R<(\ell-1)/3$.  For every $C$ and every finite set $\mathcal D$ of
integers there are infinitely many primes $p\equiv1\pmod{840}$,
$p\equiv R\pmod\ell$, having no certificate $(c,d)$ with $c\le C$ and
$d\in\mathcal D$.  The class contains no squares.
\item\emph{Computation (\S\ref{sec:census}).}  $\HC(p)\le64$ for all
$18\,803\,924\,340$ primes $p\equiv1\pmod4$ below $10^{12}$; the value
$64$ occurs only at $p=21\,841$.
\end{enumerate}

\subsection*{What is not proved}
We do not prove the Erd\H{o}s--Straus conjecture, nor $\HC(p)<\infty$ for
all large $p$, nor any bound $\HC(p)\le F(p)$ for all large $p$.  We do not
prove \eqref{eq:P415}: Theorems~Q and~E are conditional and Theorem~U
leaves a bounded number of prime factors from $\cQ$ in $4p+1$.  Theorem~W
limits one construction only; it does not bound $\HC$.  Theorem~N excludes
fixed finite certificate families only; arguments whose moduli grow with
$p$ are not excluded.  The computation is a finite verification
(\S\ref{sec:census} describes the checks); the Erd\H{o}s--Straus
conjecture itself is verified much further \cite{Salez,MD}.

\subsection*{Relation to the literature}
The solutions in Proposition~\ref{prop:slots}(c) are of Type~I or Type~II
in the sense of \cite{ET}; for instance the words with $4\mid t$ are the
case $c=1$ of the characterisation $4abcd=a+b+nc$ of Type~II solvability
in \cite[Prop.~2.7]{ET}.  The divisor criterion of Lemma~\ref{lem:lines}
for $c\equiv0\pmod 4$ is L\'opez's Type~A condition
$p\equiv-4d\pmod{4dn-1}$ and its level analogue \cite{Lopez}.  Theorem~E
adapts the method of Friedlander--Iwaniec \cite[Thm.~4, \S6]{FI09}, who
treat $\sum_{n\le x}\Lambda(n)r(n-2)r(n+2)$ under a level-of-distribution
hypothesis.  Theorem~U adapts the vector sieve of Nath and Xie
\cite{NX}, who obtain primes $p=m^2+n^2+1$ with $\Omega(p+2)$ bounded.  We
did not find the statements of Theorems~Q, E, U or N in the literature; our
search was limited to arXiv and web-indexed sources (MathSciNet and
zbMATH were not searched).  A related one-parameter statistic,
the least $a\equiv3\pmod4$ for which $(p+a)/4$ is a denominator of a
solution, is studied in the companion note \cite{Win}.

\subsection*{Conventions}
$p$ and $q$ denote primes, $\ell$ usually denotes a prime, $(\cdot/\cdot)$
is the Jacobi symbol, and $\Omega_\cQ(n)$ is the number of prime factors of
$n$ in $\cQ$, counted with multiplicity.  Constants called absolute do not
depend on any parameter.

\section{Words, certificates and the statistic \texorpdfstring{$\HC$}{HC}}\label{sec:words}

\begin{lemma}\label{lem:word}
Let $n\in\Z$ and $t,u,v\in\N$, and put $A=tu-1$, $A^*=tv-1$.
\begin{enumerate}
\item[(a)] The following are equivalent: $tuv-u-v=n$; $vA=n+u$;
$uA^*=n+v$; $AA^*=tn+1$.
\item[(b)] If $n\ge1$ and $(t,u,v)$ is a word of $n$, then $A,A^*\ge1$ and
$t,u,v\le n+2$.
\end{enumerate}
\end{lemma}

\begin{proof}
(a) $vA-u=tuv-u-v$, symmetrically for $uA^*$, and
$AA^*=t(tuv-u-v)+1$ with $t\ne0$.  (b) If $tu=1$ then $tuv-u-v=-1$, so
$A\ge1$; likewise $A^*\ge1$.  $t=(n+u+v)/(uv)\le n+2$.  By symmetry let
$u\le v$.  Then $v(tu-1)=n+u\le n+v$, so $v(tu-2)\le n$; if $tu\ge3$ this
gives $v\le n$, and if $tu=2$ then $v=n+u\le n+2$.
\end{proof}

\begin{proposition}\label{prop:slots}
Let $n\equiv1\pmod4$, $n\ge5$, and let $(t,u,v)$ be a word of $n$, with
$A=tu-1$, $A^*=tv-1$.
\begin{enumerate}
\item[(a)] If one of $t,u,v$ lies in $\cC_n$, then $4$ divides one of the
five numbers $t,u,v,A,A^*$.
\item[(b)] Conversely, if $4$ divides one of $t,u,v,A,A^*$, then one of
$t,u,v$ lies in $\cC_n$.
\item[(c)] If $4$ divides one of $t,u,v,A,A^*$, then \eqref{eq:ES} has the
following solution:
\begin{align*}
t=4s:&\quad \frac4n=\frac1{suv}+\frac1{nsv}+\frac1{nsu};\\
u=4r:&\quad \frac4n=\frac1{rA^*}+\frac1{tnr}+\frac1{tnrA^*};\\
A=4a:&\quad \frac4n=\frac1{av}+\frac1{avA^*}+\frac1{naA^*};
\end{align*}
and symmetrically for $v=4r$ and $A^*=4a$ (exchange $u\leftrightarrow v$,
$A\leftrightarrow A^*$).
\end{enumerate}
\end{proposition}

\begin{proof}
(c) Using Lemma~\ref{lem:word}(a): for $t=4s$,
$\frac{1}{suv}+\frac{1}{nsv}+\frac{1}{nsu}=\frac{n+u+v}{nsuv}=\frac{tuv}{nsuv}=\frac4n$.
For $u=4r$, the sum is $\frac{tn+A^*+1}{tnrA^*}=\frac{t(n+v)}{tnrA^*}
=\frac{uA^*}{nrA^*}=\frac4n$, since $A^*+1=tv$.  For $A=4a$, the sum is
$\frac{nA^*+n+v}{navA^*}=\frac{ntv+v}{navA^*}=\frac{v\,AA^*}{navA^*}=\frac4n$,
since $AA^*=tn+1$.

(a) If the coordinate in $\cC_n$ is divisible by $4$ there is nothing to
prove.  Suppose $u\equiv-n\pmod8$ with $u$ odd.  Then
$vA=n+u\equiv0\pmod8$, so $v$ or $A$ is divisible by $4$.  The case of $v$
is symmetric.  If $t\equiv-n\pmod 8$ with $t$ odd, then
$AA^*=tn+1\equiv1-n^2\equiv0\pmod8$, since $n$ is odd; so $4\mid A$ or
$4\mid A^*$.

(b) We may assume that none of $t,u,v$ is divisible by $4$ and, by
symmetry, that $4\mid A=tu-1$.  Then $t$ and $u$ are odd and $n=vA-u$.  If
$8\mid A$, or if $A\equiv4\pmod8$ and $v$ is even, then $n\equiv-u\pmod8$,
so $u\in\cC_n$.  Otherwise $A\equiv4\pmod8$ and $v$ is odd, so
$n\equiv4-u\pmod 8$ and $tu\equiv5\pmod8$.  As $n\equiv1\pmod4$ we get
$u\equiv3\pmod4$.  If $u\equiv3\pmod8$ then $t\equiv7$ and $n\equiv1\pmod
8$; if $u\equiv7\pmod8$ then $t\equiv3$ and $n\equiv5\pmod8$.  In both
cases $t\equiv-n\pmod8$, so $t\in\cC_n$.
\end{proof}

So a word of $n$ has a coordinate in $\cC_n$ if and only if one of its five
quantities $t,u,v,A,A^*$ is divisible by $4$, and each such word gives an
explicit solution of \eqref{eq:ES}.  All denominators in (c) are positive.

\begin{lemma}\label{lem:lines}
Let $n\ge1$ and $c\ge2$.
\begin{enumerate}
\item[(U)] $c$ is the second (equivalently, by the symmetry
$(t,u,v)\mapsto(t,v,u)$, the third) coordinate of a word of $n$ if and
only if $n+c$ has a positive divisor $d\equiv-1\pmod c$.  Given $d$, the
word is $\bigl((d+1)/c,\ c,\ (n+c)/d\bigr)$.
\item[(T)] $c$ is the first coordinate of a word of $n$ if and only if
$cn+1$ has a positive divisor $d\equiv-1\pmod c$.  Given $d$, the word is
$\bigl(c,\ (d+1)/c,\ (e+1)/c\bigr)$ with $e=(cn+1)/d$.
\end{enumerate}
\end{lemma}

\begin{proof}
(U) If $(t,c,v)$ is a word, $d=tc-1\ge1$ divides $n+c$ by
Lemma~\ref{lem:word}(a), and $d\equiv-1\pmod c$.  Conversely $d+1$ is a
positive multiple of $c$, so $t=(d+1)/c\in\N$, and $v=(n+c)/d$ satisfies
$v(tc-1)=n+c$.  (T) If $(c,u,v)$ is a word, $d=cu-1$ divides $cn+1$ by
Lemma~\ref{lem:word}(a).  Conversely, $de=cn+1\equiv1\pmod c$ and
$d\equiv-1$ give $e\equiv-1\pmod c$, so $u=(d+1)/c$ and $v=(e+1)/c$ are
positive integers with $(cu-1)(cv-1)=cn+1$.
\end{proof}

\begin{definition}\label{def:HC}
Let $n\equiv1\pmod4$.  A \emph{certificate} for $n$ is a pair $(c,d)$ with
$c\in\cC_n$ and $d$ a positive divisor of $n+c$ or of $cn+1$ satisfying
$d\equiv-1\pmod c$; we call $d$ a \emph{certificate divisor}.  The
\emph{U-line} and the \emph{T-line} at $c$ are the integers $n+c$ and
$cn+1$.  We put $\HC(n):=\min\{c:(c,d)\text{ a certificate}\}$, with
$\min\emptyset=\infty$.
\end{definition}

Since $1,2\notin\cC_n$, Lemma~\ref{lem:lines} shows that $\HC(n)$ is the
least element of $\cC_n$ that is a coordinate of some word of $n$, and by
Proposition~\ref{prop:slots} $\HC(n)<\infty$ implies that \eqref{eq:ES}
is solvable.

\begin{remark}\label{rem:lopez}
A word with $4\mid u$ corresponds to a divisor $d=4kn'-1$ of $p+4k$, i.e.\
$p\equiv-4k\pmod{4kn'-1}$; a word with $4\mid t$, $t=4k$, $u=n'$
corresponds to $p\equiv-n'\pmod{4kn'-1}$.  These are L\'opez's Type~A and
Type~B solutions \cite[Thms.~4 and~7]{Lopez}, and his Conjecture~1 is
equivalent to: every prime has a word with a coordinate divisible by $4$.
Our $\cC_p$ also contains the odd classes $\equiv-p\pmod 8$.
\end{remark}

\begin{proposition}\label{prop:small}
Let $p\equiv1\pmod4$ be prime.
\begin{enumerate}
\item[(a)] $\HC(p)\ge3$, and $\HC(p)=3$ if and only if $p\equiv5\pmod8$;
then $(1,3,(p+3)/2)$ is a word.
\item[(b)] If $p\equiv1\pmod8$, then $\HC(p)\ge4$, and $\HC(p)=4$ if and
only if $(p+4)(4p+1)$ has a prime factor $\equiv3\pmod4$.  If $\HC(p)>4$
then $\HC(p)\ge7$.
\item[(c)] If $\HC(p)<\infty$ then $\HC(p)\le p-1$.
\end{enumerate}
\end{proposition}

\begin{proof}
(a) $1,2\notin\cC_p$ because $-p$ is odd and $\not\equiv1\pmod 8$.  $3\in\cC_p$
iff $p\equiv5\pmod8$, and then $d=2\mid p+3$ with $2\equiv-1\pmod 3$.
(b) For $p\equiv1\pmod 8$, $\cC_p=4\N\cup(7+8\Z_{\ge0})$, so its two least
elements are $4$ and $7$.  The numbers $p+4$ and $4p+1$ are odd; an odd
number has a divisor $\equiv3\pmod4$ iff it has a prime factor
$\equiv3\pmod4$, because products of numbers $\equiv1\pmod4$ are
$\equiv1\pmod4$.  (c) By Lemma~\ref{lem:word}(b) every coordinate is
$\le p+2$, and $p,p+1,p+2\notin\cC_p$: they are $\equiv1,2,3\pmod4$, and
the odd ones are not $\equiv-p\pmod8$ since $2p,\,2p+2\not\equiv0\pmod8$.
\end{proof}

\subsection{Reciprocity}
The following lemma is the only arithmetic property of $p$ beyond
congruences that enters the covering results; it is the certificate form
of the classical obstruction at squares.

\begin{lemma}\label{lem:recip}
Let $n\equiv1\pmod4$ be a positive integer and $(c,d)$ a certificate for
$n$.  Then $\gcd(d,n)=1$ and $(d/n)=-1$.
\end{lemma}

\begin{proof}
Since $d\equiv-1\pmod c$, $\gcd(c,d)=1$.  Write $d=2^ad'$ with $d'$ odd.
If $d\mid n+c$ then $n\equiv-c\pmod{d'}$; if $d\mid cn+1$ then
$n\equiv-\bar c\pmod{d'}$ with $c\bar c\equiv1\pmod{d'}$.  As
$(\bar c/d')=(c/d')$, in both cases $(n/d')=(-c/d')$, which is $\pm1$.  By
Jacobi reciprocity and $n\equiv1\pmod4$, $(d'/n)=(n/d')$; this is nonzero,
so $\gcd(n,d')=1$, and $\gcd(n,2)=1$.  Hence
\begin{equation}\label{eq:recip1}
\Bigl(\frac dn\Bigr)=\Bigl(\frac2n\Bigr)^a\Bigl(\frac{-c}{d'}\Bigr).
\end{equation}

\emph{Case $4\mid c$.}  Then $d\equiv-1\pmod 4$, so $a=0$, $d'=d\equiv3\pmod4$
and $(-1/d)=-1$.  Write $c=4k$, $k=2^bk'$ with $k'$ odd.  If $b\ge1$ then
$8\mid c\mid d+1$, so $d\equiv7\pmod8$ and $(2/d)=1$.  By reciprocity, as
$(d-1)/2$ is odd, $(k'/d)=(d/k')(-1)^{(k'-1)/2}$, and $d\equiv-1\pmod{k'}$
gives $(d/k')=(-1/k')=(-1)^{(k'-1)/2}$; so $(k'/d)=1$.  Hence $(c/d)=1$
and $(-c/d)=-1$, and \eqref{eq:recip1} gives $(d/n)=-1$.

\emph{Case $c$ odd.}  Then $c\equiv-n\pmod 8$, so $c\equiv3\pmod 4$.  By
reciprocity $(c/d')=(d'/c)(-1)^{(d'-1)/2}=(d'/c)(-1/d')$, so
$(-c/d')=(d'/c)$.  From $2^ad'\equiv-1\pmod c$ we get
$(2/c)^a(d'/c)=(-1/c)=-1$.  Since $c\equiv-n\pmod 8$ and $(2/m)$ depends
only on $\pm m\bmod 8$, $(2/c)=(2/n)$.  So \eqref{eq:recip1} gives
$(d/n)=(2/n)^a\cdot(-(2/n)^a)=-1$.
\end{proof}

For a prime $p$ the lemma says that every certificate divisor is a
quadratic non-residue modulo $p$; in particular it has a prime factor that
is a non-residue.

\begin{corollary}\label{cor:squares}
\begin{enumerate}
\item[(a)] If $m$ is odd, then $m^2$ has no certificate: $\HC(m^2)=\infty$.
\item[(b)] Let $m\ge1$ and $n\equiv1\pmod{8m}$.  Then no certificate
divisor of $n$ has all its prime factors among the primes dividing $2m$.
\item[(c)] For every finite set $\mathcal D\subset\N$ there are infinitely
many primes $p$ having no certificate $(c,d)$ with $d\in\mathcal D$.
\end{enumerate}
\end{corollary}

\begin{proof}
(a) $(d/m^2)=(d/m)^2\in\{0,1\}$.  (b) $(2/n)=1$ as $n\equiv1\pmod8$, and
for an odd prime $\ell\mid m$, $(\ell/n)=(n/\ell)=(1/\ell)=1$.  So such a
$d$ has $(d/n)=1$, contrary to Lemma~\ref{lem:recip}.  (c) Apply (b) with
$m$ the product of the elements of $\mathcal D$, and Dirichlet's theorem.
\end{proof}

Part (a) is the certificate form of the vanishing of Type~I and Type~II
solutions at odd squares \cite[Prop.~1.6]{ET}.  It implies that any proof
that $\HC(p)<\infty$ must use a property of $p$ that fails at odd squares.

%%CONTINUE%%
\end{document}
